% section: 指数対数

  \section{指数対数}
    \subsection{指数}
      $a$が$p$回掛けられている状況を考える.
      ここで
      \[
        a\times a\times\cdots\times a
      \]
      のように記すのは煩雑なので,
      \[
        a\times a\times\cdots\times a=a^p
      \]
      と表す.
      このときの$p$を指数といい,
      \[
        f(x)=a^x
      \]
      で表されるような関数を指数関数という.
    \subsection{対数}
      $a^p=q$のとき,\,$p$は
      \[
        \log_a{q}=p
      \]
      がのように表される.
      この時$\log$をログと読み,\,$p$を対数という.
      また,\,$a$を底,\,$q$を真数と言う.
      一般に真数は0より大きい.
      また,
      \[
        e=\lim_{x \to \infty} \left(1+\frac{1}{x}\right)^x
      \]
      で定義される数$e$を自然対数の底またはネイピア数という.また,
      \[
        f(x)=\log_a x
      \]
      というような関数を対数関数といい,指数関数の逆関数である.
      逆関数とは
      \[
        y=f(x)
      \]
      を満たすような$x,y$について,
      \[
        x=f(y) \iff y=g(x)
      \]
      となるような関数$g$を指す.
      関数$g$はよく,\,$f^{-1}$と書かれる.

      対数に関して,一般に以下が成り立つ.
      \begin{align*}
        &\log_a{p^q}=q\log_a{p}\\
        &\log_a{MN}=\log_a{M}+\log_a{N}\\
        &\log_a{\frac{M}{N}}=\log_a{M}-\log_a{N}\\
        &\log_a{b}=\frac{\log_c{b}}{\log_c{a}}
      \end{align*}
    
